
\begin{table}[htbp]
    \centering
    \caption{Scope and Limits of the Ontological Account of Identity}
    \label{tab:identity-scope-and-limits}

    \renewcommand{\arraystretch}{1.4}
    \begin{tabular}{p{3.3cm} p{5.0cm} p{5.3cm}}
        \hline
        \textbf{Domain} &
        \textbf{Ontological Commitment} &
        \textbf{Limit or Unresolved Question} \\
        \hline

        Being and identity &
        Being and identity are ontologically inseparable. An entity does not first exist as an indeterminate substrate and subsequently acquire identity as an additional property. &
        Their analytical distinction must not be mistaken for an ontological separation or a temporal sequence between being and identity. \\

        Identity and determination &
        Identity belongs to the determinate constitution of an entity. Determination is therefore relevant to the constitution of identity. &
        Not every determination is necessarily constitutive of identity. The criteria distinguishing constitutive from non-constitutive determinations require further articulation. \\

        Identity and unity &
        Identity and unity are related dimensions of an entity's ontological constitution. Unity concerns coherence, whereas identity concerns what makes an entity the entity that it is. &
        The precise conditions under which unity contributes to identity, and whether one entity may involve multiple identities, remain open questions. \\

        Identity and individuation &
        Individuation concerns the differentiation of an entity as this entity rather than another. It is related to, but not identical with, identity. &
        The relationship between the constitution of identity and the conditions of individuation requires further ontological analysis. \\

        Identity and numerical identity &
        Numerical identity concerns the sameness of an entity across different descriptions, conditions, or temporal positions. &
        Numerical sameness alone does not provide an exhaustive account of the entity's ontological identity. \\

        Identity and invariance &
        Invariance may be relevant to the preservation or analysis of identity across transformations. &
        Identity cannot be defined in advance as invariance under a specified class of allowable transformations. The ontological status and scope of invariants require independent justification. \\

        Identity and signature &
        Signature provides a structured characterization of features through which an entity may be distinguished or described. &
        Signature cannot be presumed to exhaust, constitute, or replace the identity of the entity. \\

        Identity and composition &
        The composition of an entity may contribute to its constitution and organization. &
        Identity cannot be reduced to the sum, arrangement, or persistence of constituent elements without an account of the ontological unity of the entity. \\

        Identity and structure &
        Structural determinations may constrain and contribute to the constitution of an entity. &
        A structural description does not necessarily capture every constitutive aspect of identity. The relation between structure and identity must not be assumed to be exhaustive. \\

        Identity and expression &
        Expression concerns the articulation or expression of an entity's determinations. &
        Whether expression has an independent ontological status and how it relates to identity require separate argument. Expression must not be treated as a complete substitute for identity. \\

        Identity and manifestation &
        Manifestation concerns the ontological realization of an entity or its determinations within a particular domain. &
        Manifestation must not be reduced in advance to the appearance or representation of identity. Its ontological relation to identity and to the receiving domain requires further analysis. \\

        Identity and persistence &
        Persistence concerns the continued existence of an entity across conditions or processes. &
        Persistence does not automatically establish identity preservation. The conditions under which an entity remains itself through change remain a central ontological problem. \\

        Identity and plurality &
        A single entity may involve multiple constituents, determinations, relations, or manifestations. &
        Whether such multiplicity can correspond to multiple identities within one entity remains unresolved. Neither plurality nor unity alone settles this question. \\

        Identity and change &
        An entity may undergo changes in its determinations, relations, structure, or manifestations. &
        Neither every change entails identity loss nor every continuity guarantees identity preservation. The relevant conditions must be established through further argument. \\

        Ontological scope &
        The present account concerns identity as an aspect of being independently of any knowing subject, cognitive reconstruction, or descriptive practice. &
        Epistemic access to identity, cognitive reconstruction, and the role of the knowing subject belong to a distinct epistemological investigation and are not presupposed here. \\

        \hline
    \end{tabular}
\end{table}
